\documentclass{article}
\usepackage{graphicx} % Required for inserting images

\title{Master Thesis}
\author{Max Widmann}
\date{June 2026}

\begin{document}

\maketitle

\section{Introduction}
\section{Results}
\subsection{Experimental/Analysis Set Up}

\textbf{Imaging Set Up}\\
Our two-photon-microscope setup allows us to project a visual stimulus from underneath alive larvae fixed with agarose while % add figure of set up
simultaneously imaging fluorescence in the brain at single cell resolution at a rate of  1 to 2$Hz$. The pan-neuronally expressed calcium indicator GCaMP changes its fluorescence based on the calcium concentration in the cell which serves as an indirect read-out of the activity for each neuron. \\
The main stimuli presented for this work started with flickering
white % check this
dots which then either started to move horizontally, or stayed flickering stationary while the background turn brighter on one and darker on the other side of the midline of the fish, creating a split-view luminance stimulus.\\
%
% here could also be the part about WARP stimulus
%
As shown previously % cite armin, flo
the moving dots activate a circuit of functionally and anatomically defined neurons in the anterior hindbrain, which elicit the OMR. The luminance stimuli also activate neurons with specific reactions to the overall brightness or the change in it. These functional cell types also have validated through behavior experiments. \\% Katja
For fish aged  3, 4 and 5 dpf multiple different z-planes were imaged, focusing on the anterior hindbrain, the optic tectum and parts of the pretectum. On each of the z-planes the stimuli were presented in random order for 8 to 20 times. Some fish were imaged overnight, recording cell activity from the same z planes about every  3 hours. \\
\\
\textbf{Imaging Data Processing}\\
Using the two-photon preprocessing pipeline the fluorescence image stacks were motion-corrected, aligned to the stimulus presentations, the outlines of singles cells segmented and the average stacks registered to the reference brain ZBrain. Yielding an fluorescence time series for each segmented cell area, stimulus and trial, as well as coordinates and assigned region in the reference brain.\\ %add maybe total number of cells
Cells fluorescence traces were normalized to $\Delta F/F_{0} $ using the median of the pre-stimulus period as $F_{0}$.
%look at this again to get the correct thresholds and stuff
After this cells were selected as active if they % ...
% also add number of active cells
\\
To summarize the trials of one stimulus into one representative trial per active cell, the trial traces where clustered into two clusters using tslearn's \textit{DTW} algorithm %source, also figure with artificial and real data
with the underlying assumption that there will be one \textit{true response} and one \textit{noise} cluster. The reason for this instead of simple averaging was the observation that cells will not response to every presentation of the stimulus, which would lead to a loss of responsive signals were they averaged with the non-responsive trials. Furthermore this provided another metric to analyze stimulus responses of cells, named \textit{reliability}, which is the fraction of trials which were classified to the \textit{true response} cluster. The final trial-adjusted response trace was created by weighing the trial traces according to their distance from the \textit{true response} cluster's centroid trace.\\
After the trial-adjustment of response traces only cells with a \textit{reliability} of over 0.3 in at least half of the stimuli were used for further analysis.\\
Comparing the fraction of active cells selected by the analysis above already shows a difference between 3, 4 and 5 dpf old fish. While % ... add fractions numbers and plot


\subsection{Age dependent neuron response to luminance and motion stimuli}

As described above, the functional cell types responding most prominently and thought to be responsible for the larva's behavior, were already investigated and described previously. Therefore, it was possible to use ideal responses of these cell types as regressors to classify cells.
Two subsets of calcium response regressors were used to classify neurons from larvae of 3, 4 and 5dpf.
The first set of regressors contained the characteristic responses of hindbrain OMR  %? maybe list the cell types or put a figure I guess
cells (\textit{Motion Integrator}, \textit{Motion Onset} and \textit{Slow Motion Integrator})to motion in their preferred direction (equal to the hemisphere they are on). % refer to Armin and Flo for the regressors
The second set was used to detect cells that respond to the changes in luminance, while not responding to motion (\textit{Bright}, \textit{Dark}, \textit{Change} and \textit{Luminance Integrator}).\\ % refere to Katja for regressors
Pearson correlation of the cells trial-adjusted calcium trace to the regressor subsets was used to assign a functional cell type.
For the OMR regressor set only cells assigned to the \textit{Rhombencephalon} region of the  reference brain atlas were used, as the cell types of interest are proposed to be located there, while luminance sensitive could be located in any region of the brain.\\
This cell type classification revealed that, while characteristic luminance sensitive cells could already be found at similar frequency and with similar correlation values ($\rho$) %add numbers of cells, also figure of histograms
at the age of 3dpf then at later ages, the cell types thought to be responsible for OMR were barely detected in the hindbrain of 3dpf fish .\\ %figure
Further evidence for the later development of motion sensitive cells is that the average calcium activity of all cells located in the brain regions which had the largest change in the fraction of active cells (number of active cells divided by the total number of segmented cell assigned to that region) has a high correlation a GCaMP-kernel-convolved step function which is the proposed activity of a hindbrain motion integrator cell. \\ %this will also be shown in a figure which the regions and the average trace




\subsection{something something gene expression}

\section{Methods}
\section{Outlook}

\end{document}
